\begin{theorem}\label{fam:src:D.1.6}\label{fam:st:D.1.6}
Fix any admissible $c$. For each integral alternating matrix $E$, let
$Z_E(c)$ be the zero set on the marking cover of the explicit holomorphic
function in \eqref{fam:eq:cd-NS}, with $\beta=\beta^{\rm part}+c$ given by
\eqref{fam:eq:cd-beta}. Let $\pi$ be the cover map to $B^\circ$, and define
\[
 \mathcal E_c=\pi\!\left(\bigcup_{E\ne0,\ \operatorname{Pf}(E)=0}Z_E(c)\right),
 \qquad
 \mathcal A_c=\pi\!\left(\bigcup_{\operatorname{Pf}(E)<0}Z_E(c)\right).
\]
These give the exact loci of all elliptic quotients and all abelian
surfaces, respectively. The algebraic dimension of every smooth fibre is
\[
 a(F_t)=
 \begin{cases}
 2,&t\in\mathcal A_c,\\
 1,&t\in\mathcal E_c\setminus\mathcal A_c,\\
 0,&t\notin\mathcal E_c\cup\mathcal A_c.
 \end{cases}
\]
Together with \eqref{fam:eq:cd-NS}, these are an exact classification for every
$t$, by countable unions of explicitly given analytic zero sets; they do
not assert that the loci have a finite algebraic description. They include
all quotient maps, without requiring preservation of the first coordinate.
\end{theorem}
\Uses{\cref{fam:eq:cd-NS,fam:src:D.1.4}; \cite[Theorem 4.1, p.~1, and Theorem 4.4, p.~3]{Schnell19}}
\begin{proof}
\label{app:rederive:fam-src-D.1.6}
We apply the Appell--Humbert correspondence and the Riemann projectivity
criterion in \cite[Theorem 4.1, p.~1, and Theorem 4.4, p.~3]{Schnell19}; the reduction to the
present explicit periods is as follows. The real vector-space isomorphism
$\Pi:\mathbb R^4\to\mathbb C^2$ transfers $E$ to the actual alternating
form $\widetilde E(v,w)=E(\Pi^{-1}v,\Pi^{-1}w)$ on the tangent space.
Equation \eqref{fam:eq:cd-NS} gives $J$-invariance. The explicit comparison is
\[
 H_E(v,w)=\widetilde E(v,Jw)-i\widetilde E(v,w).
\]
Direct substitution gives $H_E(Jv,w)=iH_E(v,w)$ and
$H_E(w,v)=\overline{H_E(v,w)}$. Its associated real Chern form is
$-\operatorname{Im}H_E=\widetilde E$, since
$(i/2)(H_E(v,w)-H_E(w,v))=\widetilde E(v,w)$.
Thus the Appell--Humbert integrality condition is exactly integrality
of the given $E$, with its sign fully specified. The complex-oriented square computed in \ref{fam:src:D.1.4} shows that
$\operatorname{Pf}(E)<0$ is equivalent to definiteness of this Hermitian
form. Its sign can be chosen positive. The cited projectivity criterion
then gives precisely the asserted abelian locus, in both directions.

For nonzero $E$ with Pfaffian zero, its real rank is two. Its nullspace
$K_{\mathbb R}\subset\mathbb R^4$ is a rational plane because it is the
kernel of an integral matrix. $J$-invariance makes $\Pi K_{\mathbb R}$ a
complex line, and $K=K_{\mathbb R}\cap\mathbb Z^4$ is saturated. The maps
\[
 (\Pi K_{\mathbb R})/\Pi K\ \hookrightarrow F_t
 \ \longrightarrow\
 \mathbb C^2/(\Pi K_{\mathbb R}+\Pi\mathbb Z^4)
\]
are holomorphic group morphisms; the first is a closed embedded elliptic
curve and the second is its quotient onto an elliptic curve, with
connected kernel. A basis of the saturated $K$ extends to an integral basis of $\mathbb Z^4$.
The other two periods project to real-linearly independent vectors in the
one-dimensional complex quotient, because $\Pi$ is a real isomorphism.
They therefore form a discrete lattice there. Compactness follows by
taking the image of its compact real fundamental parallelogram. Conversely any elliptic quotient
pulls an integral positive area form back to a nonzero rank-two integral
$(1,1)$ form. This proves the characterization of $\mathcal E_c$.
Equivalently all its kernels can be found by saturated rank-two lattices
with a basis $v,w$ satisfying
$\det_{\mathbb C}(\Pi v,\Pi w)=0$ and real independence. This is a second
exact integer test, with explicit inclusion and quotient morphisms.

For the algebraic-dimension assertion, an effective divisor on a torus
has a semipositive Hermitian form: averaging its integration current under
translations gives that constant form. A nonzero semipositive form has
rank one or two. Rank two gives the abelian case. Rank one gives the
elliptic quotient just constructed. In the absence of a definite form,
there cannot be two different such kernel subtori. Indeed addition from
their product onto $F_t$ would be a surjective holomorphic homomorphism
with finite kernel, and the product polarization would descend after
multiplication by an integer, giving a definite form. Thus there is a
single kernel $K_t$ in this case. Every effective divisor is a sum of
its translates: its intersection number with $K_t$ is zero, whereas a
curve mapping nonconstantly to the elliptic quotient has strictly positive
intersection with a general quotient fibre. Every meromorphic function
therefore factors through this quotient. One can check the last statement
on a fibre avoiding its polar divisor, where it is a holomorphic function
on a compact elliptic curve and hence constant; analytic continuation
then gives the factorization. The quotient supplies nonconstant
meromorphic functions, so the algebraic dimension is exactly one.
If neither locus occurs, a nonconstant meromorphic function would have
a polar divisor of one of the two forbidden ranks. This proves dimension
zero. Finally monodromy acts by integral determinant-one changes of basis,
preserving \eqref{fam:eq:cd-NS} and its Pfaffian. Thus the stated unions descend
without any marking ambiguity. Each individual matrix and every locus
is specified; no expected identification is left in these formulas.
\end{proof}
